\tinysidebar{\debug{exercises/{59-27-pointers-and-arrays/question.tex}}}Implement the pointer version of bubblesort and binary search:



\verb!int * bubblesort(int * start, * end)!



\verb!int * binarysearch(int * start, * end, int target)!



In your \texttt{main()}, get integers \texttt{n}, \texttt{start}, \texttt{end}, \texttt{target} from the user and then do the following:



\begin{itemize}

\item

  Allocate memory for a dynamic integer array of size \texttt{n}. and put

  random integers (from 0..9) into the array.

\item

  Print the array.

\item

  Bubblesort \texttt{x[start]}, \ldots, \texttt{x[end-1]}.

\item

  Print the array.

\item

  Perform binary search for \texttt{target} in \texttt{x[start]},

  \ldots, \texttt{x[end-1]} and print the address and the value at

  that address.

\end{itemize}



Solution is below. Don't memorize the code. Understand the motivation and intention of bubblesort and binarysearch, implement them using index values, then convert to pointers.



\begin{consolethree}[escapeinside=||]

#include <iostream>



// print *start, .., *(end - 1)

void println(int * start, int * end) // or begin and end

{

    for (int * p = start; p < end; ++p)

    {

        std::cout << (*p) << ' ';

    }

    std::cout << '\n';

}



// bubblesort *start, .., *(end - 1)

void bubblesort(int * start, int * end)

{

    for (int * q = end - 2; q >= start; --q)

    {

        for (int * p = start; p <= q; ++p)

        {

            if (*p > *(p + 1))

            {

                int t = *p; *p = *(p + 1); *(p + 1) = t;

            }

        }

    }

}



// binarysearch for target in *start, .., *(end - 1)

int * binarysearch(int * start, int * end, int target)

{

    while (start < end)

    {

        int * mid = start + (end - start) / 2;

        if (*mid == target)

        {

            return mid;

        }

        else if (*mid > target)

        {

            end = mid;

        }

        else

        {

            start = mid + 1;

        }

    }

    return NULL;

}



int main()

{

    int n, start, end, target;

    std::cin >> n >> start >> end >> target;

    int * x = new int[n];

    for (int i = 0; i < n; ++i)

    {

        x[i] = rand() % 10;

    }

    print(&x[0], &x[n]);

    bubblesort(&x[start], &x[end]);

    print(&x[0], &x[n]);

    int * p = binarysearch(&x[start], &x[end], target);

    std::cout << p << ' '

              << (p != NULL ? *p : -1) << '\n';

    return 0;

}

\end{consolethree}

